For fatigue monitoring, the way each baseline failed in the strategy comparison is more informative than the size of the margin.
The thresholding approach used in the BLINKER-threshold baseline~\citep{kleifges2017blinker} retained most blinks but at the expense of lower precision. In a drowsiness-monitoring application, such false detections could artificially inflate blink-rate indices used to assess fatigue~\citep{tian2021fatigue}.
This limitation also helps explain why the complete BLINKER pipeline incorporates additional morphological fitting and channel-selection steps rather than relying on thresholding alone.
MNE~\citep{GramfortEtAl2013a} instead missed most blinks in several sessions on each dataset. This is arguably the more serious failure for a warning system, which must be reliable for each driver rather than on average~\citep{ingre2006subjective}.
%
Both failure patterns are consistent with calibrating the threshold on the whole recording.
In the BLINKER-threshold baseline, blink-free epochs may dilute the pool and shift the threshold towards the background EEG, so smaller non-blink deflections could cross it.
In MNE, a single large artefact raises the range-based threshold above many genuine blinks.
The first effect is a particular concern in driving, because the degree of dilution depends on how often the drivs blinks, and blink frequency changes with sleepiness and differ considerably between drivers~\citep{schleicher2008blinks}.
A threshold calibrated in this way would therefore depend partly on the state it is meant to monitor.
The proposed pipeline is designed to avoid this dilution by estimating its threshold only from epochs flagged as containing prominent amplitude excursions.
The proposed approach outperformed BLINKER by a similar F1 margin on Internal and Cao2018, even though the two datasets were recorded with different hardware and montages.
Its advantage therefore does not appear to depend on a particular recording system, although both datasets were collected during simulated lane-keeping tasks.